%% =============================================================================
%%  Revision de las observaciones recibidas el 07-11-2026: evaluacion de su estructura y de su procedencia.
%%  Documento aparte del articulo. Compilar con pdfLaTeX. La tabla, los conteos y la lista de observaciones que no tienen sentido
%%  los escribe codigo/gen_revision_observaciones.py desde codigo/observaciones_revision.py y desde las cifras del manuscrito.
%% =============================================================================
\documentclass[10pt,a4paper]{article}
\usepackage[utf8]{inputenc}
\usepackage[T1]{fontenc}
\usepackage[spanish,es-tabla,es-nodecimaldot]{babel}
\usepackage[landscape,a4paper,margin=1.7cm]{geometry}
\usepackage{amsmath,amssymb}
\usepackage{booktabs}
\usepackage{longtable}
\usepackage{array}
\usepackage[table]{xcolor}
\usepackage{enumitem}
\usepackage[colorlinks=true,linkcolor=blue!50!black,urlcolor=blue!50!black]{hyperref}
\definecolor{okverde}{rgb}{0.84,0.94,0.84}
\definecolor{okamarillo}{rgb}{1.0,0.95,0.72}
\definecolor{okrojo}{rgb}{0.98,0.80,0.78}
\definecolor{okgris}{rgb}{0.86,0.88,0.92}
\setlength{\parskip}{4pt}
\setlength{\parindent}{0pt}
\setlist{itemsep=2pt,topsep=3pt}
\emergencystretch=2em
% Generado por gen_tablas_opt.py (no editar a mano)
\newcommand{\DelayOfFijo}{13.81}
\newcommand{\DelayOfAct}{10.89}
\newcommand{\DelayOfActAf}{9.63}
\newcommand{\DelayOfActV}{8.34}
\newcommand{\DelayOfMax}{17.44}
\newcommand{\DelayOfPap}{10.95}
\newcommand{\DelayOfPapAf}{9.67}
\newcommand{\DelayOfMdp}{9.35}
\newcommand{\DelayOfFlc}{7.91}
\newcommand{\DelayOfFlcC}{9.97}
\newcommand{\DelayOfMlp}{8.09}
\newcommand{\DelayOfMlpC}{9.46}
\newcommand{\DelayNvFijo}{13.00}
\newcommand{\DelayNvAct}{10.55}
\newcommand{\DelayNvActAf}{9.17}
\newcommand{\DelayNvActV}{7.82}
\newcommand{\DelayNvMax}{16.55}
\newcommand{\DelayNvPap}{10.40}
\newcommand{\DelayNvPapAf}{9.30}
\newcommand{\DelayNvMdp}{8.77}
\newcommand{\DelayNvFlc}{7.40}
\newcommand{\DelayNvFlcC}{9.37}
\newcommand{\DelayNvMlp}{7.53}
\newcommand{\DelayNvMlpC}{8.95}
\newcommand{\DelayVaFijo}{12.29}
\newcommand{\DelayVaAct}{9.23}
\newcommand{\DelayVaActAf}{8.07}
\newcommand{\DelayVaActV}{6.65}
\newcommand{\DelayVaMax}{19.00}
\newcommand{\DelayVaPap}{8.65}
\newcommand{\DelayVaPapAf}{8.06}
\newcommand{\DelayVaMdp}{7.94}
\newcommand{\DelayVaFlc}{6.07}
\newcommand{\DelayVaFlcC}{8.41}
\newcommand{\DelayVaMlp}{6.24}
\newcommand{\DelayVaMlpC}{8.05}
\newcommand{\DelaySaFijo}{40.67}
\newcommand{\DelaySaAct}{54.06}
\newcommand{\DelaySaActAf}{24.13}
\newcommand{\DelaySaActV}{25.61}
\newcommand{\DelaySaMax}{152.25}
\newcommand{\DelaySaPap}{105.87}
\newcommand{\DelaySaPapAf}{58.11}
\newcommand{\DelaySaFlc}{24.29}
\newcommand{\DelaySaFlcC}{25.74}
\newcommand{\DelaySaMlp}{25.10}
\newcommand{\DelaySaMlpC}{24.29}
\newcommand{\RedFijo}{$-$26.1}
\newcommand{\RedAct}{0.6}
\newcommand{\RedActAf}{12.1}
\newcommand{\RedActV}{23.9}
\newcommand{\RedMax}{$-$59.3}
\newcommand{\RedPapAf}{11.7}
\newcommand{\RedMdp}{14.7}
\newcommand{\RedFlc}{27.8}
\newcommand{\RedFlcC}{9.0}
\newcommand{\RedMlp}{26.2}
\newcommand{\RedMlpC}{13.6}
\newcommand{\RedPapOf}{27.8}
\newcommand{\RedActAfOf}{17.8}
\newcommand{\RedActVOf}{5.1}
\newcommand{\RedVistaPapOf}{23.9}
\newcommand{\RedVistaAfOf}{13.4}
\newcommand{\RedPapNv}{28.8}
\newcommand{\RedActAfNv}{19.3}
\newcommand{\RedActVNv}{5.3}
\newcommand{\RedVistaPapNv}{24.8}
\newcommand{\RedVistaAfNv}{14.8}
\newcommand{\RedPapVa}{29.8}
\newcommand{\RedActAfVa}{24.7}
\newcommand{\RedActVVa}{8.7}
\newcommand{\RedVistaPapVa}{23.2}
\newcommand{\RedVistaAfVa}{17.6}
\newcommand{\RedPapSa}{77.1}
\newcommand{\RedActAfSa}{$-$0.7}
\newcommand{\RedActVSa}{5.2}
\newcommand{\RedVistaPapSa}{75.8}
\newcommand{\RedVistaAfSa}{$-$6.2}
\newcommand{\ParMediaPap}{28.2}
\newcommand{\ParMejorPap}{16}
\newcommand{\ParPeorPap}{0}
\newcommand{\ParMinPap}{19.5}
\newcommand{\ParMaxPap}{49.3}
\newcommand{\ParMediaActAf}{22.0}
\newcommand{\ParMejorActAf}{16}
\newcommand{\ParPeorActAf}{0}
\newcommand{\ParMinActAf}{2.9}
\newcommand{\ParMaxActAf}{46.5}
\newcommand{\ParMediaActV}{7.7}
\newcommand{\ParMejorActV}{14}
\newcommand{\ParPeorActV}{0}
\newcommand{\ParMinActV}{0.5}
\newcommand{\ParMaxActV}{26.0}
\newcommand{\ParMediaAct}{31.3}
\newcommand{\ParMejorAct}{16}
\newcommand{\ParPeorAct}{0}
\newcommand{\ParMinAct}{11.4}
\newcommand{\ParMaxAct}{56.8}
\newcommand{\ParMediaFijo}{44.1}
\newcommand{\ParMejorFijo}{16}
\newcommand{\ParPeorFijo}{0}
\newcommand{\ParMinFijo}{39.2}
\newcommand{\ParMaxFijo}{61.2}
\newcommand{\ParMediaMdp}{19.2}
\newcommand{\ParMejorMdp}{16}
\newcommand{\ParPeorMdp}{0}
\newcommand{\ParMinMdp}{2.7}
\newcommand{\ParMaxMdp}{43.5}
\newcommand{\DelayOfFijos}{8.44}
\newcommand{\DelayNvFijos}{7.88}
\newcommand{\DelayVaFijos}{6.56}
\newcommand{\DelaySaFijos}{24.80}
\newcommand{\RedActFijos}{12.3}
\newcommand{\RedPapFijos}{22.9}
\newcommand{\DelayOfReglas}{8.47}
\newcommand{\DelayNvReglas}{7.93}
\newcommand{\DelayVaReglas}{6.55}
\newcommand{\DelaySaReglas}{25.00}
\newcommand{\RedActReglas}{12.0}
\newcommand{\RedPapReglas}{22.6}
\newcommand{\DelayOfGSiete}{8.60}
\newcommand{\DelayNvGSiete}{8.14}
\newcommand{\DelayVaGSiete}{6.85}
\newcommand{\DelaySaGSiete}{25.15}
\newcommand{\RedActGSiete}{10.6}
\newcommand{\RedPapGSiete}{21.4}
\newcommand{\DelayOfGCinco}{8.21}
\newcommand{\DelayNvGCinco}{7.75}
\newcommand{\DelayVaGCinco}{6.41}
\newcommand{\DelaySaGCinco}{25.16}
\newcommand{\RedActGCinco}{14.8}
\newcommand{\RedPapGCinco}{25.1}
\newcommand{\DelayOfGCuatro}{8.04}
\newcommand{\DelayNvGCuatro}{7.58}
\newcommand{\DelayVaGCuatro}{6.20}
\newcommand{\DelaySaGCuatro}{24.38}
\newcommand{\RedActGCuatro}{16.5}
\newcommand{\RedPapGCuatro}{26.6}
\newcommand{\DelayOfSemCero}{8.00}
\newcommand{\DelayNvSemCero}{7.51}
\newcommand{\DelayVaSemCero}{6.21}
\newcommand{\DelaySaSemCero}{24.79}
\newcommand{\RedActSemCero}{16.9}
\newcommand{\RedPapSemCero}{26.9}
\newcommand{\DelayOfSemUno}{8.04}
\newcommand{\DelayNvSemUno}{7.49}
\newcommand{\DelayVaSemUno}{6.24}
\newcommand{\DelaySaSemUno}{24.63}
\newcommand{\RedActSemUno}{16.5}
\newcommand{\RedPapSemUno}{26.6}
\newcommand{\DelayOfSemDos}{7.91}
\newcommand{\DelayNvSemDos}{7.40}
\newcommand{\DelayVaSemDos}{6.07}
\newcommand{\DelaySaSemDos}{24.29}
\newcommand{\RedActSemDos}{17.8}
\newcommand{\RedPapSemDos}{27.8}
\newcommand{\DelayOfLinC}{9.53}
\newcommand{\DelayNvLinC}{8.97}
\newcommand{\DelayVaLinC}{8.10}
\newcommand{\DelaySaLinC}{25.11}
\newcommand{\RedActLinC}{1.0}
\newcommand{\RedPapLinC}{13.0}
\newcommand{\RedActFijo}{$-$43.4}
\newcommand{\RedActAct}{$-$13.1}
\newcommand{\RedActActV}{13.4}
\newcommand{\RedActMax}{$-$81.2}
\newcommand{\RedActPap}{$-$13.8}
\newcommand{\RedActPapAf}{$-$0.4}
\newcommand{\RedActMdp}{2.9}
\newcommand{\RedActFlc}{17.8}
\newcommand{\RedActFlcC}{$-$3.5}
\newcommand{\RedActMlp}{16.0}
\newcommand{\RedActMlpC}{1.7}
\newcommand{\ResDifOf}{0.2}
\newcommand{\ResDifOfMax}{2.1}
\newcommand{\ResDifSa}{4.1}
\newcommand{\ResDifSaMax}{13.9}
\newcommand{\EscNNv}{64}
\newcommand{\EscPapNvRed}{29.7}
\newcommand{\EscPapNvLo}{27.6}
\newcommand{\EscPapNvHi}{31.8}
\newcommand{\EscPapNvP}{$<\!0.001$}
\newcommand{\EscPapNvMejor}{64}
\newcommand{\EscPapNvPeor}{0}
\newcommand{\EscActAfNvRed}{23.3}
\newcommand{\EscActAfNvLo}{20.1}
\newcommand{\EscActAfNvHi}{26.5}
\newcommand{\EscActAfNvP}{$<\!0.001$}
\newcommand{\EscActAfNvMejor}{64}
\newcommand{\EscActAfNvPeor}{0}
\newcommand{\EscActVNvRed}{7.9}
\newcommand{\EscActVNvLo}{6.0}
\newcommand{\EscActVNvHi}{9.8}
\newcommand{\EscActVNvP}{$<\!0.001$}
\newcommand{\EscActVNvMejor}{53}
\newcommand{\EscActVNvPeor}{11}
\newcommand{\EscNVa}{64}
\newcommand{\EscPapVaRed}{30.0}
\newcommand{\EscPapVaLo}{28.3}
\newcommand{\EscPapVaHi}{31.6}
\newcommand{\EscPapVaP}{$<\!0.001$}
\newcommand{\EscPapVaMejor}{64}
\newcommand{\EscPapVaPeor}{0}
\newcommand{\EscActAfVaRed}{27.1}
\newcommand{\EscActAfVaLo}{24.8}
\newcommand{\EscActAfVaHi}{29.3}
\newcommand{\EscActAfVaP}{$<\!0.001$}
\newcommand{\EscActAfVaMejor}{64}
\newcommand{\EscActAfVaPeor}{0}
\newcommand{\EscActVVaRed}{10.1}
\newcommand{\EscActVVaLo}{8.8}
\newcommand{\EscActVVaHi}{11.5}
\newcommand{\EscActVVaP}{$<\!0.001$}
\newcommand{\EscActVVaMejor}{64}
\newcommand{\EscActVVaPeor}{0}
\newcommand{\EscNSa}{32}
\newcommand{\EscPapSaRed}{69.5}
\newcommand{\EscPapSaLo}{64.3}
\newcommand{\EscPapSaHi}{74.7}
\newcommand{\EscPapSaP}{$<\!0.001$}
\newcommand{\EscPapSaMejor}{32}
\newcommand{\EscPapSaPeor}{0}
\newcommand{\EscActAfSaRed}{$-$0.3}
\newcommand{\EscActAfSaLo}{$-$1.3}
\newcommand{\EscActAfSaHi}{0.6}
\newcommand{\EscActAfSaP}{0.36}
\newcommand{\EscActAfSaMejor}{14}
\newcommand{\EscActAfSaPeor}{18}
\newcommand{\EscActVSaRed}{4.4}
\newcommand{\EscActVSaLo}{3.0}
\newcommand{\EscActVSaHi}{5.9}
\newcommand{\EscActVSaP}{$<\!0.001$}
\newcommand{\EscActVSaMejor}{30}
\newcommand{\EscActVSaPeor}{2}
\newcommand{\MedOfFijos}{8.47}
\newcommand{\SdOfFijos}{0.04}
\newcommand{\MedNvFijos}{7.89}
\newcommand{\SdNvFijos}{0.03}
\newcommand{\MedVaFijos}{6.56}
\newcommand{\SdVaFijos}{0.02}
\newcommand{\MedSaFijos}{24.94}
\newcommand{\SdSaFijos}{0.13}
\newcommand{\NSemFijos}{3}
\newcommand{\RedPapMedFijos}{22.7}
\newcommand{\RedActAfMedFijos}{12.0}
\newcommand{\RedActVMedFijos}{$-$1.6}
\newcommand{\MedOfReglas}{8.44}
\newcommand{\SdOfReglas}{0.03}
\newcommand{\MedNvReglas}{7.87}
\newcommand{\SdNvReglas}{0.05}
\newcommand{\MedVaReglas}{6.55}
\newcommand{\SdVaReglas}{0.00}
\newcommand{\MedSaReglas}{24.90}
\newcommand{\SdSaReglas}{0.11}
\newcommand{\NSemReglas}{3}
\newcommand{\RedPapMedReglas}{23.0}
\newcommand{\RedActAfMedReglas}{12.4}
\newcommand{\RedActVMedReglas}{$-$1.2}
\newcommand{\MedOfGSiete}{8.55}
\newcommand{\SdOfGSiete}{0.05}
\newcommand{\MedNvGSiete}{8.11}
\newcommand{\SdNvGSiete}{0.04}
\newcommand{\MedVaGSiete}{6.84}
\newcommand{\SdVaGSiete}{0.02}
\newcommand{\MedSaGSiete}{24.70}
\newcommand{\SdSaGSiete}{0.40}
\newcommand{\NSemGSiete}{3}
\newcommand{\RedPapMedGSiete}{21.9}
\newcommand{\RedActAfMedGSiete}{11.2}
\newcommand{\RedActVMedGSiete}{$-$2.5}
\newcommand{\MedOfGCinco}{8.19}
\newcommand{\SdOfGCinco}{0.05}
\newcommand{\MedNvGCinco}{7.72}
\newcommand{\SdNvGCinco}{0.08}
\newcommand{\MedVaGCinco}{6.44}
\newcommand{\SdVaGCinco}{0.09}
\newcommand{\MedSaGCinco}{25.15}
\newcommand{\SdSaGCinco}{0.09}
\newcommand{\NSemGCinco}{3}
\newcommand{\RedPapMedGCinco}{25.2}
\newcommand{\RedActAfMedGCinco}{14.9}
\newcommand{\RedActVMedGCinco}{1.7}
\newcommand{\MedOfGCuatro}{8.05}
\newcommand{\SdOfGCuatro}{0.02}
\newcommand{\MedNvGCuatro}{7.56}
\newcommand{\SdNvGCuatro}{0.01}
\newcommand{\MedVaGCuatro}{6.20}
\newcommand{\SdVaGCuatro}{0.04}
\newcommand{\MedSaGCuatro}{24.50}
\newcommand{\SdSaGCuatro}{0.11}
\newcommand{\NSemGCuatro}{3}
\newcommand{\RedPapMedGCuatro}{26.5}
\newcommand{\RedActAfMedGCuatro}{16.4}
\newcommand{\RedActVMedGCuatro}{3.4}
\newcommand{\MedOfGTres}{7.98}
\newcommand{\SdOfGTres}{0.06}
\newcommand{\MedNvGTres}{7.47}
\newcommand{\SdNvGTres}{0.06}
\newcommand{\MedVaGTres}{6.18}
\newcommand{\SdVaGTres}{0.09}
\newcommand{\MedSaGTres}{24.57}
\newcommand{\SdSaGTres}{0.25}
\newcommand{\NSemGTres}{3}
\newcommand{\RedPapMedGTres}{27.1}
\newcommand{\RedActAfMedGTres}{17.1}
\newcommand{\RedActVMedGTres}{4.2}
\newcommand{\PdMaxDif}{0.7}

% Generado por gen_revision_observaciones.py (no editar a mano)
\newcommand{\RevN}{21}
\newcommand{\RevNProcede}{11}
\newcommand{\RevNParte}{4}
\newcommand{\RevNNo}{5}
\newcommand{\RevNNoVer}{1}
\newcommand{\RevNLevantadas}{15}

\newcommand{\FLC}{FLC-O}
\newcommand{\PMR}{Pappis y Mamdani reconstruido}
\newcommand{\alerta}[1]{\begin{center}\fcolorbox{red!60!black}{red!4}{\parbox{0.96\textwidth}{#1}}\end{center}}
\newcommand{\caja}[1]{\begin{center}\fcolorbox{blue!40!black}{blue!4}{\parbox{0.96\textwidth}{#1}}\end{center}}

\begin{document}

\begin{center}
{\Large\bfseries Revisión de las observaciones recibidas el 07-11-2026}\\[3pt]
{\large Evaluación de su estructura y de su procedencia, y qué se hizo con cada una}\\[4pt]
{\small Artículo «Mejora de un controlador de semáforos de lógica difusa para una intersección aislada: optimización del controlador de Pappis y Mamdani (1977)»}\\[2pt]
{\small Raul Nolberto Quispe Saya y Kevin Luis Valero Ccacya}
\end{center}

% =============================================================================
\section{Resumen}

\caja{\textbf{Qué se recibió.} Una revisión de 12 páginas (\texttt{Revision\_Valero}) con formato de plantilla: una etapa de verificación de lo entregado, puntuaciones, nueve criterios técnicos (C1 a C9), una matriz de reproducibilidad, cinco observaciones mayores, diez menores, una evaluación de datos, de figuras y de notación, y la recomendación «Major Rev».\\[3pt]
\textbf{Cómo está estructurada.} Bien en lo esencial: cada observación mayor trae problema, severidad, impacto técnico y arreglo, y casi todas dicen en qué sección o tabla están. Tiene, sin embargo, errores propios que conviene conocer (Sec.~\ref{sec:estructura}).\\[3pt]
\textbf{Qué se concluyó.} De las \RevN{} observaciones distintas que se extrajeron del documento, \RevNProcede{} proceden y se levantaron; \RevNParte{} proceden en parte (se levantó lo que tiene fondo); \RevNNo{} \textbf{no proceden} (son artefactos de extracción de texto, referencias cruzadas erróneas, algo ya cumplido o un error factual del revisor) y \RevNNoVer{} no es verificable por el revisor. Todo lo que procede está corregido en la nueva versión del artículo; el cuadro de lo levantado es un documento aparte.}

\alerta{\textbf{Observaciones que no tienen sentido, para tenerlas presentes.}\\[2pt]
\begin{itemize}
\item \textbf{O6 (NO PROCEDE).} Pie de la Tabla 3: «CONTROLDORES» debe ser «CONTROLADORES». \emph{Por qu\'e no tiene sentido:} No existe tal error. El fuente dice «Controladores comparados y la información que usan» y el texto extraído de nuestro PDF muestra «C ONTROLADORES COMPARADOS»: el título está en versalitas y la extracción de texto pierde letras. El propio revisor declara ruido de extracción en esas páginas.
\item \textbf{O14 (NO PROCEDE).} Tabla S2 del suplemento: agregar unidades a las columnas «Error medio» y «Error máximo». \emph{Por qu\'e no tiene sentido:} Las unidades ya están: los encabezados dicen «Error medio ( \%)» y «Error máximo ( \%)».
\item \textbf{O15 (NO PROCEDE).} Justificación, Sec. 7: reemplazar «venta» por «ventana». \emph{Por qu\'e no tiene sentido:} No aparece «venta» en el texto: el fuente y el PDF dicen «ventana». El revisor declara ruido de extracción en la página 10 de ese documento.
\item \textbf{O19 (NO VERIFICABLE).} El repositorio no se pudo verificar externamente (sin acceso web); estado «UNVERIFIABLE». \emph{Por qu\'e no tiene sentido:} Es una limitación del entorno del revisor, no un defecto del manuscrito. Se comprobó que el repositorio es público (responde sin credenciales) y que un clon limpio de main pasa todas las pruebas y compila el artículo.
\item \textbf{O20 (NO PROCEDE).} Tabla 3: «no hay ecuaciones explícitas de las pertenencias»; la lista de nudos de la Tabla 3 no basta para reconstruir. \emph{Por qu\'e no tiene sentido:} Referencia cruzada errónea: la Tabla 3 es la de controladores comparados; los nudos están en el pie de la Tabla 2. La sustancia (ecuaciones exactas) es válida y se atiende en O2.
\item \textbf{O21 (NO PROCEDE).} Arreglo propuesto en C3: «hacer de la comparación con FLC-O causal (mejora de \RedActMlpC{} \% sobre el actuado) la comparación justa principal». \emph{Por qu\'e no tiene sentido:} Error factual del revisor. El \RedActMlpC{} \% es de la política neuronal causal; FLC-O causal cambia el retardo en \RedActFlcC{} \% respecto del actuado afinado, es decir, es peor (\DelayOfFlcC{} contra \DelayOfActAf{} s/veh). Además, la comparación causal no es la principal para el resultado de 27.8 \%, pues el controlador de 1977 sí tiene vista adelante. La alternativa que el revisor ofrece (actuado con vista adelante) sí se hizo (O1).
\end{itemize}
}

\textbf{La observación que más importa.} La asimetría de información (O1) era cierta y cambió el resultado. Con la misma vista adelante, la ventaja de \FLC{} sobre el actuado afinado pasa de \RedActAfOf\,\% a \RedActVOf\,\% en el conjunto oficial. Lo que no cambió es la comparación principal, contra el controlador de 1977 reconstruido (\RedFlc\,\%), porque ese controlador ya ve las llegadas de los próximos 10\,s. El detalle está en la Sec.~\ref{sec:hallazgo}.

% =============================================================================
\section{Evaluación de la estructura del documento recibido}
\label{sec:estructura}

\textbf{Lo que está bien estructurado.}
\begin{itemize}
\item Las cinco observaciones mayores siguen el formato problema, severidad, impacto técnico y arreglo accionable, y cada una cita ubicaciones concretas del manuscrito.
\item Separa lo que es defecto del manuscrito de lo que es limitación de la extracción de texto (la sección «Extraction Integrity Report») y declara que no hay deficiencias bloqueantes.
\item Distingue observaciones mayores de menores y las agrupa por tema (reproducibilidad, estadística, presentación).
\item Reconoce fortalezas reales: separación entrenamiento, validación y prueba, estadística pareada con Holm y verificación exacta del simulador.
\end{itemize}

\textbf{Lo que no está bien estructurado o es incorrecto.}
\begin{itemize}
\item \textbf{Duplicación.} La falta de ecuaciones exactas de las pertenencias aparece al menos ocho veces (Etapa 0, C4, C5, C9, matriz de reproducibilidad, observación mayor 2, Fig.~2 y Tabla~3); la falta de dispersión, otras seis. Contarlas por separado infla la lista: aquí se agrupan en \RevN{} observaciones distintas.
\item \textbf{Referencias cruzadas erróneas.} Dice que los nudos están en la «Tabla 3», pero están en el pie de la Tabla~2 (la Tabla~3 es la de controladores comparados); ubica las referencias [25] y [26] en la Sec.~IV-B, pero están en la Sec.~II; habla de un «entrenamiento de 20 minutos» que no existe en ningún documento (el manuscrito decía 15 y el tiempo medido es de 13) y lo ubica en la Sec.~V-A, que trata del simulador.
\item \textbf{Artefactos de extracción presentados como defectos.} Pide corregir «CONTROLDORES» y «venta», y agregar unidades que ya estaban. El propio documento reconoce ruido de extracción en esas páginas, y los tres puntos son consistentes con haber trabajado sobre el texto extraído y no sobre el PDF compilado.
\item \textbf{Un arreglo con un error factual.} Propone tomar como comparación justa la de «FLC-O causal (mejora de 1.7 \% sobre el actuado)», pero ese 1.7 \% es de la política neuronal causal; FLC-O causal es peor que el actuado (\RedActFlcC\,\%).
\item \textbf{Una observación sobre algo que el revisor no pudo comprobar.} El estado «no verificable» del repositorio es una limitación de su entorno, no del manuscrito.
\item \textbf{Observaciones ya cumplidas.} Pide aclarar que la banda de la Fig.~3 es solo de \FLC{}, cuando el pie ya lo decía; y pide dispersión en la Tabla~4 cuando el suplemento ya la traía para tres controladores.
\item \textbf{Puntuaciones sin criterio.} Las notas numéricas (novedad 3/5, reproducibilidad 3/5, etc.) no traen la regla con que se asignan y no son accionables; no se cuentan como observaciones.
\end{itemize}

% =============================================================================
\section{El hallazgo principal: la información igualada}
\label{sec:hallazgo}

Para responder a la observación O1 se implementó un actuado afinado que recibe la misma información que el controlador de 1977 y que \FLC{}: las llegadas de los próximos 10\,s. Se afinó con los mismos escenarios de entrenamiento y el mismo rango de parámetros que el resto (\texttt{codigo/opt\_afinar\_referencias.py}, que además reproduce exactamente el afinado anterior del actuado y de la presión máxima).

\begin{center}
\small
\begin{tabular}{@{}lrrrrrr@{}}
\toprule
Conjunto de prueba & \PMR{} & Actuado afinado & Actuado con vista adelante & \FLC{} & \FLC{} frente al actuado & \FLC{} frente al actuado con vista \\
\midrule
Oficial (16 pares) & \DelayOfPap{} & \DelayOfActAf{} & \DelayOfActV{} & \DelayOfFlc{} & $-$\RedActAfOf\,\% & $-$\RedActVOf\,\% \\
No vistos (64) & \DelayNvPap{} & \DelayNvActAf{} & \DelayNvActV{} & \DelayNvFlc{} & $-$\RedActAfNv\,\% & $-$\RedActVNv\,\% \\
Variable (64) & \DelayVaPap{} & \DelayVaActAf{} & \DelayVaActV{} & \DelayVaFlc{} & $-$\RedActAfVa\,\% & $-$\RedActVVa\,\% \\
Saturado (32) & \DelaySaPap{} & \DelaySaActAf{} & \DelaySaActV{} & \DelaySaFlc{} & \RedActAfSa\,\% & $-$\RedActVSa\,\% \\
\bottomrule
\end{tabular}

{\footnotesize\textit{Fuente:} \texttt{resultados/opt/evaluacion\_final.json}; retardo medio en s/veh; la reducción es sobre el promedio del conjunto.}
\end{center}

Lo que implica, dicho sin suavizar:
\begin{itemize}
\item La comparación principal no cambia: \FLC{} reduce \RedFlc\,\% el retardo de \PMR{}, que tiene la misma información, y es significativamente mejor en \ParMejorPap{} de 16 pares.
\item La ventaja sobre el actuado era en gran parte información: el actuado afinado gana \RedVistaAfOf\,\% solo con recibir la vista adelante. Con información igualada, \FLC{} mantiene una ventaja de \RedActVOf\,\%, significativa en \ParMejorActV{} de 16 pares y sin ningún par peor, pero es mucho menor que la que se decía.
\item Un actuado simple con la misma vista adelante ya mejora en \RedVistaPapOf\,\% al controlador de 1977. Una parte grande de la distancia al controlador de 1977 se recorre sin lógica difusa.
\item Cerca de la saturación, \FLC{} y el actuado afinado siguen empatados (sin diferencia significativa, $p=\EscActAfSaP$), y el actuado con vista adelante es algo peor (\DelaySaActV{}).
\end{itemize}
El manuscrito lo dice así en el resumen, la Sec.~VIII-A, la discusión y las conclusiones.

% =============================================================================
\section{Evaluación observación por observación}
\label{sec:tabla}

Veredictos: \colorbox{okverde}{PROCEDE} se corrigió; \colorbox{okamarillo}{PROCEDE EN PARTE} tiene un fondo válido que se corrigió y una parte que no; \colorbox{okrojo}{\textbf{NO PROCEDE}} no se corrigió porque no hay defecto o está mal formulada; \colorbox{okgris}{NO VERIFICABLE} depende del entorno del revisor.

\begingroup\footnotesize\setlength{\tabcolsep}{4pt}
\begin{longtable}{@{}p{7.4cm}p{2.1cm}p{8.4cm}p{7.3cm}@{}}
\caption{Evaluaci\'on de cada observaci\'on recibida}\label{tab:evaluacion}\\
\toprule
Observaci\'on recibida & Veredicto & Evaluaci\'on (raz\'on y evidencia) & Qu\'e se hizo en el manuscrito \\
\midrule
\endfirsthead
\toprule
Observaci\'on recibida & Veredicto & Evaluaci\'on (raz\'on y evidencia) & Qu\'e se hizo en el manuscrito \\
\midrule
\endhead
\textbf{O1}\ (Mayor 1): Asimetría de información: FLC-O usa 10 s de vista adelante y el actuado afinado no; la ventaja de \RedActAfOf{} \% sobre el actuado es en parte una ventaja de información. Arreglo propuesto: agregar un actuado con vista adelante, o presentar la comparación sin vista adelante como la comparación justa y calificar el 17.8 \%.\par\smallskip{\itshape Ubicaci\'on citada: Abstract; Sec. VIII-A; Tabla 4; Sec. IX; C3; Sec. 7} & \cellcolor{okverde}\textbf{PROCEDE} & Es cierto y era la objeción más importante. Se implementó y afinó un actuado con la misma vista adelante (4 parámetros, mismos escenarios de entrenamiento y rango). Resultado: la ventaja de FLC-O sobre el actuado cae de \RedActAfOf{} \% a \RedActVOf{} \% en el conjunto oficial (significativa en \ParMejorActV{} de 16 pares, \ParPeorActV{} peor), a \RedActVNv{} \% con caudales no vistos, a \RedActVVa{} \% con demanda variable y a \RedActVSa{} \% cerca de la saturación. Matiz que el revisor no vio: la comparación principal, con el controlador de 1977 reconstruido, ya tenía información igual (ese controlador también ve las llegadas de los próximos 10 s). & Nuevo rival «actuado afinado con vista adelante» en el protocolo, en las Tablas 3 a 6, en la Fig. 3 y en el resumen; el \RedActAfOf{} \% se califica como «sin vista adelante» y se reporta junto al \RedActVOf{} \%. Se agrega, con honestidad, que el actuado con vista adelante ya mejora \RedVistaPapOf{} \% al controlador de 1977.\par\smallskip{\itshape Queda en: Sec. VII (p. 6); Tabla 3 (p. 6); Sec. VIII-A (p. 7); Tabla 4 (p. 8); Tabla 5 (p. 9); Fig. 3 (p. 8); Sec. VIII-B (p. 7); Tabla 6 (p. 10); Sec. IX (p. 9); Sec. III del suplemento (p. 1)} \\
\midrule
\textbf{O2}\ (Mayor 2): Las ecuaciones exactas de las funciones de pertenencia de FLC-O (formas, hombros, comportamiento fuera del universo) no están en el manuscrito; la Fig. 2 y la lista de nudos no bastan para reconstruirlo sin el código.\par\smallskip{\itshape Ubicaci\'on citada: Fig. 2; Tabla 3; Justificación Sec. 7; C4, C5, C9; Parte 4B ítem 12} & \cellcolor{okverde}\textbf{PROCEDE} & Era cierto: el texto describía la partición en palabras y daba los nudos con un decimal. Se verificó la corrección con una implementación independiente escrita solo con las ecuaciones nuevas. & Se agregan las ecuaciones de las tres pertenencias, la parametrización de los nudos, el comportamiento fuera de [k1, k3]; los nudos y pesos con 3 decimales; en el suplemento, los 42 parámetros con 6 decimales y un ejemplo numérico. Nueva prueba test\_flco\_ecuaciones.py: la implementación independiente coincide con el código (diferencia 4e-16) y con la tabla impresa.\par\smallskip{\itshape Queda en: Sec. VI-B (p. 5); Ec. (6) (p. 5); Fig. 2 (p. 7); Tabla 2 (p. 6); Sec. IX del suplemento (p. 2); Tabla S7 del suplemento (p. 5); Tabla S8 del suplemento (p. 5)} \\
\midrule
\textbf{O3}\ (Mayor 3): Los conjuntos de prueba no oficiales (Tabla 6) y las filas de ablación (Tabla 7) carecen de número de réplicas y de dispersión; la ablación usa una sola semilla de entrenamiento. Arreglo: agregar n, semillas y SD/IC a las Tablas 6 y 7; estadística multi-semilla para verde mínimo, número de reglas y pesos de la vista adelante.\par\smallskip{\itshape Ubicaci\'on citada: Tablas 6 y 7; C5, C6; Sec. 7} & \cellcolor{okverde}\textbf{PROCEDE} & Cierto. En la Tabla 6 cada escenario es un par de caudales con una réplica, pero eso no estaba dicho; y la ablación comparaba entrenamientos únicos, sin saber cuánto varía una semilla de otra. & Tabla 6: n, semillas, media con IC 95 \% y reducción por escenario con IC y prueba de Wilcoxon. Tabla 7: media y desviación entre 3 semillas de entrenamiento para verde mínimo (3, 4, 5 y 7 s), horizontes fijos y 81 reglas; las filas con una sola semilla se marcan como exploratorias (n = 1). Se entrenaron 10 modelos más para esto.\par\smallskip{\itshape Queda en: Sec. VIII-B (p. 7); Tabla 6 (p. 10); Sec. VIII-C (p. 7); Tabla 7 (p. 10); Tabla S6 del suplemento (p. 4)} \\
\midrule
\textbf{O4}\ (Mayor 4): La comparación con Pappis y Mamdani es contra una reconstrucción calibrada, no contra la tabla impresa; el resumen y las etiquetas de las tablas no lo califican de forma consistente. Arreglo: etiquetar «reconstruido con A/Q calibradas» y dar el error de 9.2 \% junto a la comparación principal.\par\smallskip{\itshape Ubicaci\'on citada: Sec. IV-B; Sec. IX; Tabla 4; Abstract} & \cellcolor{okverde}\textbf{PROCEDE} & Cierto en la presentación: el manuscrito lo declaraba en la Sec. IV-B y en la Sec. IX, pero el resumen y las tablas decían «Pappis y Mamdani» a secas. Aclaración: las reglas, los conjuntos y los operadores son los impresos; solo las entradas A y Q son una definición de este trabajo. & Se usa «Pappis y Mamdani reconstruido» en el resumen, el protocolo, las tablas y la discusión; el error de 9.2 \% (máximo 18.6 \%) acompaña la definición en el protocolo y en el resumen.\par\smallskip{\itshape Queda en: Sec. IV-B (p. 3); Sec. VII (p. 6); Tabla 3 (p. 6); Tabla 4 (p. 8); Sec. IX (p. 9)} \\
\midrule
\textbf{O5}\ (Mayor 5): La Tabla 4 solo da retardos medios; la Fig. 3(a) tiene banda de confianza solo para FLC-O. Arreglo: columnas de SD o IC 95 \% en la Tabla 4 y bandas para todos los controladores.\par\smallskip{\itshape Ubicaci\'on citada: Tabla 4; Fig. 3; C7} & \cellcolor{okverde}\textbf{PROCEDE} & Cierto en el texto principal. La dispersión por par ya estaba en el suplemento (Tabla S5) para tres controladores, pero no en la Tabla 4. & Tabla 4: media ± desviación estándar de las 20 réplicas para cada controlador. Fig. 3(a): banda de IC 95 \% para todos (la cota es determinista). Suplemento S5: se agrega el actuado con vista adelante.\par\smallskip{\itshape Queda en: Tabla 4 (p. 8); Fig. 3 (p. 8); Tabla S5 del suplemento (p. 4)} \\
\midrule
\textbf{O6}\ (Menor): Pie de la Tabla 3: «CONTROLDORES» debe ser «CONTROLADORES».\par\smallskip{\itshape Ubicaci\'on citada: Tabla 3, pie} & \cellcolor{okrojo}\textbf{NO PROCEDE} & No existe tal error. El fuente dice «Controladores comparados y la información que usan» y el texto extraído de nuestro PDF muestra «C ONTROLADORES COMPARADOS»: el título está en versalitas y la extracción de texto pierde letras. El propio revisor declara ruido de extracción en esas páginas. & Ninguna (no hay defecto en el manuscrito).\par\smallskip{\itshape Queda en: Tabla 3 (p. 6)} \\
\midrule
\textbf{O7}\ (Menor): Pie de la Fig. 3: aclarar que la banda es el IC 95 \% solo de FLC-O; agregar IC para los otros o decirlo explícitamente.\par\smallskip{\itshape Ubicaci\'on citada: Fig. 3, pie} & \cellcolor{okamarillo}\textbf{PROCEDE EN PARTE} & El pie ya decía «la banda es el intervalo de confianza de 95 \% de FLC-O». Se acepta la sugerencia de fondo (bandas para todos). & Bandas para todos los controladores y pie reescrito (ver O5).\par\smallskip{\itshape Queda en: Fig. 3 (p. 8)} \\
\midrule
\textbf{O8}\ (Menor): Pie de la Tabla 6: indicar número de réplicas y semillas de cada conjunto de prueba.\par\smallskip{\itshape Ubicaci\'on citada: Tabla 6, pie} & \cellcolor{okverde}\textbf{PROCEDE} & Cierto: faltaba. & El pie y los encabezados dan el número de unidades (16 pares, 64, 64 y 32 escenarios) y las semillas de cada conjunto.\par\smallskip{\itshape Queda en: Tabla 6 (p. 10); Sec. VII (p. 6)} \\
\midrule
\textbf{O9}\ (Menor): Sec. IX: «Pappis y Mamdani» debe ser «Pappis y Mamdani reconstruido» en la primera mención del 27.8 \%.\par\smallskip{\itshape Ubicaci\'on citada: Sec. IX} & \cellcolor{okverde}\textbf{PROCEDE} & Cierto (ver O4). & Cambiado en la discusión, las conclusiones y el resumen.\par\smallskip{\itshape Queda en: Sec. IX (p. 9); Sec. X (p. 9)} \\
\midrule
\textbf{O10}\ (Menor): Resumen: calificar el 17.8 \% sobre el actuado como «con vista adelante de 10 s».\par\smallskip{\itshape Ubicaci\'on citada: Abstract} & \cellcolor{okverde}\textbf{PROCEDE} & Cierto, y la revisión de fondo (O1) muestra que es la parte más sensible del resultado. & El resumen en los dos idiomas separa «actuado sin vista adelante: \RedActAfOf{} \%» de «actuado con la misma vista adelante: \RedActVOf{} \%».\par\smallskip{\itshape Queda en: Sec. VIII-A (p. 7)} \\
\midrule
\textbf{O11}\ (Menor): Sec. IV-B: revisar la consistencia de las referencias [25] «Vuong et al. (ANFIS)» y [26] «Shaheen et al.» con la lista de referencias.\par\smallskip{\itshape Ubicaci\'on citada: Sec. IV-B} & \cellcolor{okamarillo}\textbf{PROCEDE EN PARTE} & Se verificaron ambas contra Crossref (título, autores, revista, volumen, páginas, año): son consistentes. La ubicación citada es errónea: las dos citas están en la Sec. II (estado del arte), no en la IV-B. Solo faltaba el número de artículo de Shaheen et al. & Se completa la entrada de Shaheen et al. (número 15273).\par\smallskip{\itshape Queda en: Sec. II (p. 2)} \\
\midrule
\textbf{O12}\ (Menor): Sec. V-A: indicar el modelo exacto de CPU y el tiempo de ejecución de la afirmación del «entrenamiento de 20 minutos»; la tabla de reproducibilidad marca el hardware como incompleto.\par\smallskip{\itshape Ubicaci\'on citada: Sec. V-A; Parte 4B ítem 7} & \cellcolor{okamarillo}\textbf{PROCEDE EN PARTE} & El dato de hardware sí faltaba. Pero no existe una afirmación de 20 minutos en ningún documento: el manuscrito decía 15 y el tiempo medido es 787 s (13 min). Además, el tiempo de entrenamiento está en la Sec. VI-D, no en la V-A (que trata del simulador). El «hardware» de este renglón de la plantilla es solo la computadora de las simulaciones: el estudio no tiene equipos físicos, instrumentación ni programas de PLC. & Se agrega el procesador (Intel Xeon a 2.10 GHz, 4 vCPU, 15 GB) en las Sec. V-A y VI-D y en la disponibilidad, y el tiempo medido (13 min, 787 s). Se declara además que es el único equipo del estudio, que es solo de simulación (introducción, «Alcance»).\par\smallskip{\itshape Queda en: Sec. V-A (p. 5); Sec. VI-D (p. 6); «Disponibilidad del código» (p. 9); Sec. I (p. 1)} \\
\midrule
\textbf{O13}\ (Menor): Tabla 5: aclarar si el intervalo es sobre las 16 mejoras por par o un IC de la media.\par\smallskip{\itshape Ubicaci\'on citada: Tabla 5} & \cellcolor{okverde}\textbf{PROCEDE} & Era el rango (mínimo a máximo) de las 16 mejoras, no un intervalo de confianza; el rótulo «Intervalo» confundía. & Columna renombrada «Rango», se agrega «IC 95 \%» de la media y el pie lo explica.\par\smallskip{\itshape Queda en: Tabla 5 (p. 9)} \\
\midrule
\textbf{O14}\ (Menor): Tabla S2 del suplemento: agregar unidades a las columnas «Error medio» y «Error máximo».\par\smallskip{\itshape Ubicaci\'on citada: Suplemento, Tabla S2} & \cellcolor{okrojo}\textbf{NO PROCEDE} & Las unidades ya están: los encabezados dicen «Error medio ( \%)» y «Error máximo ( \%)». & Ninguna.\par\smallskip{\itshape Queda en: Tabla S2 del suplemento (p. 3)} \\
\midrule
\textbf{O15}\ (Menor): Justificación, Sec. 7: reemplazar «venta» por «ventana».\par\smallskip{\itshape Ubicaci\'on citada: Justificación, Sec. 7} & \cellcolor{okrojo}\textbf{NO PROCEDE} & No aparece «venta» en el texto: el fuente y el PDF dicen «ventana». El revisor declara ruido de extracción en la página 10 de ese documento. & Ninguna.\par\smallskip{\itshape Queda en: sin cambios} \\
\midrule
\textbf{O16}\ (mayor, clasif. propia): Novedad incremental: el resumen y la introducción deben declarar que es un estudio empírico de optimización y comparación, no un nuevo paradigma de control.\par\smallskip{\itshape Ubicaci\'on citada: C2; Scorecard; Sec. 1} & \cellcolor{okverde}\textbf{PROCEDE} & Es una calificación razonable de lo que el trabajo aporta. & Se declara en el resumen (dos idiomas) y en la introducción: contribución empírica, no una nueva teoría de control.\par\smallskip{\itshape Queda en: Sec. I (p. 1)} \\
\midrule
\textbf{O17}\ (menor, clasif. propia): Notación: definir G, R, V, Q\_G, Q\_V, u y c\_ijk en la tabla de notación antes de su primer uso.\par\smallskip{\itshape Ubicaci\'on citada: Sec. 8; Sec. VI} & \cellcolor{okamarillo}\textbf{PROCEDE EN PARTE} & G, R, V, u y c\_ijk ya estaban en la Tabla 1, antes de la Sec. VI. «Q\_G» no existe como símbolo aparte: en la Sec. IV-B era la misma cola que Q\_g. Sí faltaban Q\_V, a\_e, q\_e, los nudos k\_i y las pertenencias. & Se unifica Q\_G como Q\_g y se agregan Q\_V, a\_e, q\_e, k1-k3 y las pertenencias a la Tabla 1.\par\smallskip{\itshape Queda en: Tabla 1 (p. 4); Sec. IV-B (p. 3)} \\
\midrule
\textbf{O18}\ (mayor, clasif. propia): El «empate» cerca de la saturación (\DelaySaFlc{} frente a \DelaySaActAf{} s/veh) se reporta sin prueba de significancia.\par\smallskip{\itshape Ubicaci\'on citada: Sec. 7; Tabla 6; C-tabla de evidencia} & \cellcolor{okverde}\textbf{PROCEDE} & Cierto. Con la prueba pareada por escenario no hay diferencia significativa (p = \EscActAfSaP{}; reducción media por escenario de \EscActAfSaRed{} \%, IC 95 \% de \EscActAfSaLo{} \% a \EscActAfSaHi{} \%). Un p alto no demuestra equivalencia y así se dice. & Se reporta el IC y el p de Wilcoxon en la Tabla 6 y en la Sec. VIII-B, con la salvedad.\par\smallskip{\itshape Queda en: Sec. VIII-B (p. 7); Tabla 6 (p. 10)} \\
\midrule
\textbf{O19}\ (mayor, clasif. propia): El repositorio no se pudo verificar externamente (sin acceso web); estado «UNVERIFIABLE».\par\smallskip{\itshape Ubicaci\'on citada: Etapa 0; C9; Parte 4B ítems 10 y 11} & \cellcolor{okgris}\textbf{NO VERIFICABLE} & Es una limitación del entorno del revisor, no un defecto del manuscrito. Se comprobó que el repositorio es público (responde sin credenciales) y que un clon limpio de main pasa todas las pruebas y compila el artículo. & Para facilitar la verificación se agrega al manuscrito el commit con el que se generaron los resultados y el entorno exacto.\par\smallskip{\itshape Queda en: «Disponibilidad del código» (p. 9)} \\
\midrule
\textbf{O20}\ (menor, clasif. propia): Tabla 3: «no hay ecuaciones explícitas de las pertenencias»; la lista de nudos de la Tabla 3 no basta para reconstruir.\par\smallskip{\itshape Ubicaci\'on citada: Sec. 8 (Tabla 3)} & \cellcolor{okrojo}\textbf{NO PROCEDE} & Referencia cruzada errónea: la Tabla 3 es la de controladores comparados; los nudos están en el pie de la Tabla 2. La sustancia (ecuaciones exactas) es válida y se atiende en O2. & Ninguna adicional a O2.\par\smallskip{\itshape Queda en: Tabla 2 (p. 6)} \\
\midrule
\textbf{O21}\ (mayor, clasif. propia): Arreglo propuesto en C3: «hacer de la comparación con FLC-O causal (mejora de \RedActMlpC{} \% sobre el actuado) la comparación justa principal».\par\smallskip{\itshape Ubicaci\'on citada: C3; Sec. 5, punto 1} & \cellcolor{okrojo}\textbf{NO PROCEDE} & Error factual del revisor. El \RedActMlpC{} \% es de la política neuronal causal; FLC-O causal cambia el retardo en \RedActFlcC{} \% respecto del actuado afinado, es decir, es peor (\DelayOfFlcC{} contra \DelayOfActAf{} s/veh). Además, la comparación causal no es la principal para el resultado de 27.8 \%, pues el controlador de 1977 sí tiene vista adelante. La alternativa que el revisor ofrece (actuado con vista adelante) sí se hizo (O1). & Se mantiene la comparación sin vista adelante tal como estaba (\RedActMlpC{} \% la política neuronal; \RedActFlcC{} \% el difuso causal) y se agrega la comparación con información igualada.\par\smallskip{\itshape Queda en: Sec. VIII-C (p. 7)} \\
\midrule
\bottomrule
\end{longtable}
\endgroup


% =============================================================================
\section{Qué se hizo además de cambiar el texto}

\begin{itemize}
\item \textbf{Experimentos nuevos.} El actuado con vista adelante (afinado y evaluado en los cuatro conjuntos de prueba); diez entrenamientos más de las variantes de la ablación, para tener tres semillas en verde mínimo, horizontes fijos y 81 reglas; estadística de dispersión por unidad para todos los controladores.
\item \textbf{Reproducibilidad.} El afinado de las referencias, que antes no tenía un script, ahora es \texttt{codigo/opt\_afinar\_referencias.py} y reproduce los resultados guardados con 16 cifras. Los 42 parámetros de \FLC{} están completos en el suplemento y una prueba nueva (\texttt{codigo/test\_flco\_ecuaciones.py}) comprueba que las ecuaciones del suplemento, implementadas de forma independiente, dan la misma salida que el código.
\item \textbf{Referencias.} Se agregó la de Ruspini (verificada en Crossref) y se completó la de Shaheen y colaboradores.
\end{itemize}

\section{Límites y pendientes}
\begin{itemize}
\item El actuado con vista adelante es de una sola familia y de cuatro parámetros; otro más elaborado podría reducir la ventaja de \FLC{}. Está dicho en las amenazas a la validez.
\item Las filas de la ablación con una sola semilla (política neuronal, difuso causal, política lineal y \PMR{} afinado) siguen siendo exploratorias, y la tabla las marca como tales.
\item Con 32 escenarios saturados no se puede afirmar que \FLC{} y el actuado afinado sean equivalentes, solo que no se detecta diferencia.
\item Los tiempos de ejecución se midieron en un contenedor de 4 vCPU; en otro equipo variarán.
\end{itemize}

\end{document}
